% Reusable visual acceptance template. Override these macros before input.
\providecommand\VisualTheme{floral}
\providecommand\VisualAspect{169}
\documentclass[aspectratio=\VisualAspect]{beamer}
\usetheme{\VisualTheme}
\ifdefined\VisualBaseline
  \newcommand\florathisframe[1]{}
  \newcommand\florapalette[1]{}
  % Match the package's normal defaults for attribution without the package.
  \setbeamersize{text margin left=2.5cm,text margin right=1cm}
\else
  \usepackage{beamer-floral}
  \florasetup{seed=3,density=1.2}
\fi
\author[Floral]{The Floral Project}
\institute[Watercolour and type]{Procedural illustration and typography}
\date{}
\newcommand\VisualCase[1]{\typeout{VISUAL-CASE: #1}}
\newcommand\VisualTitle[3]{%
  \VisualCase{#1}%
  \title[Floral margins]{#2}\subtitle{#3}%
  \ifdefstring{\VisualTheme}{floral}{%
    \florathisframe{layout=left,tune={gesture={top_inset=0}}}%
  }{}%
  \begin{frame}[plain]\titlepage\end{frame}}
\begin{document}
\VisualTitle{01 short title}{Floral}{Watercolour meets typography}
\VisualTitle{02 medium title}{A garden around the words: colour, rhythm and space}{Procedural borders for thoughtful presentations}
\VisualTitle{03 long title}{Drawing botanical borders around a changing text block with procedural watercolour and carefully balanced typography}{Stems, leaves and blossoms share the page with the words}
\VisualTitle{04 very long title}{The shape of a floral presentation: composing stems, leaves, blossoms and ornaments around titles, paragraphs and figures while preserving space for the ideas within}{}
\VisualTitle{05 long title and subtitle}{Between the garden and the page: a study of procedural ornament, measured composition and the quiet spaces that bring a presentation into balance}{Following the border from the first blossom to the final flourish, with room for both a longer title and a carefully placed subtitle}
\section[Composition]{Composition}
\subsection[Typography]{Typography}
\VisualCase{06 short frame title}
\begin{frame}{Room to grow}
  \framesubtitle{The margin is part of the composition}
  \begin{itemize}\item The words define the space.\item The border grows around them.\end{itemize}
\end{frame}
\VisualCase{07 medium frame title}
\begin{frame}{Procedural watercolour borders: colour, rhythm and typographic balance}
  \framesubtitle{A longer heading leaves room for a supporting line}
  Short paragraphs let the illustration breathe.
\end{frame}
\VisualCase{08 very long frame title}
\begin{frame}{From a single stem to a complete border: arranging botanical details around the changing geometry of a presentation}
  \framesubtitle{A supporting line beneath an extended heading}
  The border gives the words a setting.
\end{frame}
\VisualCase{09 untitled after long title}
\begin{frame}
  \begin{block}{A quieter page}Without a heading, the border follows the space beneath the running head.\end{block}
\end{frame}
\subsection[Blocks]{Blocks}
\VisualCase{10 blocks and lists}
\begin{frame}{Colour and hierarchy}
  \begin{block}{Structure}A shared palette connects text and illustration.\end{block}
  \begin{alertblock}{Emphasis}A deeper ink draws attention.\end{alertblock}
  \begin{exampleblock}{Example}Foliage supplies a second accent.\end{exampleblock}
\end{frame}
\VisualCase{11 columns and table}
\begin{frame}{Balance across the page}
  \begin{columns}[T,onlytextwidth]
    \begin{column}{.46\textwidth}
      \begin{block}{Branches of an idea}
        \begin{itemize}\item A branching idea.\begin{itemize}\item A supporting detail.\begin{itemize}\item A finer point.\end{itemize}\end{itemize}\end{itemize}
      \end{block}
    \end{column}
    \begin{column}{.46\textwidth}
      \begin{tabular}{lr}Element&Weight\\\hline Leaves&3\\Blossoms&2\\Sprays&1\end{tabular}
    \end{column}
  \end{columns}
\end{frame}
\subsection[Overlays]{Overlays}
\VisualCase{12 overlays - two PDF pages}
\begin{frame}
  \frametitle{\only<1>{A steady border}\only<2>{A steady border around an unfolding idea: keeping the composition in place}}
  \only<1>{The first idea takes its place.}
  \only<2>{A second idea opens up the page.\par\vspace{1cm}The illustration stays still as the words change.}
\end{frame}
\VisualCase{13 plain content frame}
\begin{frame}[plain]{An open canvas}
  With no running head or footline, the border has the page to itself.
\end{frame}
\VisualCase{14 blue palette}
\florapalette{blue}
\begin{frame}{Blue palette}
  \begin{block}{Cool inks}Blue blossoms and foliage share an ink family.\end{block}
  \begin{alertblock}{Emphasis}Floral uses a deeper violet accent.\end{alertblock}
  \begin{exampleblock}{Continuity}Stock themes keep their own block colours.\end{exampleblock}
\end{frame}
\VisualCase{15 green palette}
\florapalette{green}
\begin{frame}{Green palette}
  \begin{block}{Living colour}Blue blossoms meet green foliage.\end{block}
  \begin{alertblock}{Emphasis}Floral uses a plum accent.\end{alertblock}
  \begin{exampleblock}{Continuity}Stock themes keep their own block colours.\end{exampleblock}
\end{frame}
\end{document}
